\section{互信息：知道一件事，能帮你猜中另一件事多少}

% ==================================================================
\subsection{先说人话：你每天早上都在用互信息}

先不讲公式，讲个你天天遇到的事。

早上 7 点，你站在宿舍门口，脑子里有两个问号：
\begin{enumerate}[leftmargin=1.8em,itemsep=1pt,topsep=2pt]
  \item \textbf{今天下雨吗？}（记为 $X$）
  \item \textbf{我要不要带伞？}（记为 $Y$）
\end{enumerate}
你对第一个问号是\emph{不确定}的：也许是 $30\%$ 会下雨，也许是 $90\%$。
这种“心里没底”的程度，信息论里有个专门的度量，叫\textbf{不确定性}
（后面你会看到它就是熵 $H(X)$）。

现在楼下大爷喊了一句：“今儿下雨！”（这是一个消息）。
这一句话，让你对“带不带伞”的犹豫几乎消失了——原来你还在纠结，
现在你基本确定要带伞。

学长要强调的是这句话的\textbf{价值}：大爷这句话给你带来的东西，
不是“下雨”这两个字，而是那句更抽象的话——
\begin{center}
\fbox{\parbox{0.86\textwidth}{\centering
\textbf{你关于“带不带伞”的不确定性，减少了多少}\\[2pt]
\footnotesize 减少的这部分不确定性，就是这句话的\textbf{信息量}。
}}
\end{center}
\textbf{互信息（mutual information）衡量的就是这个“减少了多少”。}
写成公式的样子，就是本讲的主角 $\I{X;Y}$，读作“$X$ 和 $Y$ 的互信息”，
意思是“\textbf{知道 $Y$ 之后，我对 $X$ 的不确定性减少了多少}（单位 bit）”。

\begin{idea}{本讲唯一的“一根筋”思路}
整讲你只需要记住一句话，剩下的全是这句话的变体：

\medskip
\textbf{信息量 $=$ 不确定性的减少量。}

\medskip
\begin{itemize}[leftmargin=1.4em,itemsep=1pt,topsep=1pt]
  \item 一件\textbf{本来就没准}的事（$p=0.5$），告诉你了，信息量\textbf{大}；
  \item 一件\textbf{本来就很准}的事（$p=0.99$），再说一遍，信息量\textbf{小}；
  \item 知道 $Y$ 之后 $X$ 的不确定性降低得越多，$\I{X;Y}$ 就越大。
\end{itemize}
\end{idea}

\subsubsection{先把这四个符号认全（不认全后面全是天书）}

\begin{center}
\small
\begin{tabular}{@{}lll@{}}
\toprule
符号 & 怎么读 & 大白话 \\
\midrule
$p(x)$ & “$p$ 括号 $x$” & 事件“$X$ 取值为 $x$”发生的\textbf{概率}，一个 $0\sim1$ 的数 \\
$I(x)$ & “$x$ 的自信息” & 事件 $x$ 真的发生了这件事，本身\textbf{含有多少信息量} \\
 & & $I(x)=-\lb p(x)$，单位 bit \\
$H(X)$ & “$X$ 的熵” & $X$ \textbf{整体的平均不确定性}，$H(X)=\sum_x p(x)I(x)$ \\
$H(X|Y)$ & “在 $Y$ 的条件下 $X$ 的条件熵” & \textbf{知道 $Y$ 之后}，$X$ 还剩多少不确定性 \\
$\I{X;Y}$ & “$X$ 和 $Y$ 的互信息” & 知道 $Y$ 之后，$X$ 的不确定性\textbf{减少了多少} \\
\bottomrule
\end{tabular}
\end{center}

还有一条本讲要反复用的“换算器”（考试不给你，要背）：
\[
\lb 2=1,\qquad \lb 3\approx1.58496\approx1.585,\qquad
\lb 5\approx2.32193\approx2.322,\qquad \lb 10\approx3.32193\approx3.322 .
\]
用法举例：$\lb 0.4=\lb\frac{2}{5}=\lb 2-\lb 5=1-2.32193=-1.32193$，
$\lb 0.3=\lb\frac{3}{10}=1.58496-3.32193=-1.73697$。
看到小数就用“拆成分数、再用上面三个数”的办法，别去按计算器蒙。

另外约定一件事：本讲用小写 $x,y$ 表示\textbf{一个具体的取值}（一个事件），
用大写 $X,Y$ 表示\textbf{整个随机变量}（所有取值一起）。这个大小写之别，
正是本讲前半段（事件互信息）和后半段（平均互信息）的区别。

% ==================================================================
\subsection{事件的互信息：一件事能“泄露”多少关于另一件事的信息}

\subsubsection{定义，以及分子分母到底是什么}

\begin{defbox}{定义（事件之间的互信息）}
设 $X,Y$ 是两个离散随机变量，$p(x,y)$ 读作“$x$ 和 $y$ \textbf{同时}发生的概率”
（联合概率），$p(x)$、$p(y)$ 分别是“$x$ 单独发生的概率”和“$y$ 单独发生的概率”
（边缘概率）。事件 $y$ 对事件 $x$ 的\textbf{互信息}定义为
\[
\I{x;y}=\lb\frac{p(x,y)}{p(x)\,p(y)} .
\]
底取 $2$ 时单位是 \textbf{bit}（本课程一律用底 $2$）。
\end{defbox}

\begin{idea}{三个零件的分工}
把 $\I{x;y}=\lb\dfrac{p(x,y)}{p(x)p(y)}$ 拆开看，它其实在问一句话：
\textbf{“这两件事一起发生的频率，比‘各管各的’应该有的频率高多少倍？”}
\begin{itemize}[leftmargin=1.4em,itemsep=1pt,topsep=1pt]
  \item \textbf{分母 $p(x)p(y)$}：如果 $X$ 和 $Y$ \emph{完全没关系}
        （这种情况叫\textbf{统计独立}），那么同时发生的概率就应该是
        “$p(x)$ 乘 $p(y)$”。所以分母是\textbf{“没关系时本该有的概率”}，
        也就是你的\textbf{基准线}（瞎猜的命中率）。
  \item \textbf{分子 $p(x,y)$}：这两件事\textbf{实际}同时发生的概率。
  \item \textbf{比值}：实际 $\div$ 应该 $=1$ 表示“跟没关系一模一样”；
        $>1$ 表示“一起出现得异常频繁，它俩有牵连（正相关）”；
        $<1$ 表示“它俩老是互相躲着（负相关）”。
\end{itemize}
最后套上 $\lb$，是把“倍数”翻译成“bit”，这个下文马上解释。
\end{idea}

\subsubsection{为什么一定要套对数？（这里用了一个技巧）}

因为\textbf{对数能把乘法变成加法}，而“独立”正好是乘法。

具体地：如果 $x$ 和 $y$ 独立，则 $p(x,y)=p(x)p(y)$，于是联合自信息
$I(x,y)=-\lb p(x,y)=-\lb p(x)-\lb p(y)=I(x)+I(y)$——
\textbf{两件互不相干的事，信息量正好相加}（“明天太阳升起”和“明天下雨”
这两条消息合起来的信息量，就该等于两条相加）。取对数就是为了保住这个
“独立 $\Rightarrow$ 相加”的性质；不取对数，加法就凑不出来了。

再看比特的含义：比值 $=2$（实际发生的机会是瞎猜的 $2$ 倍）
$\Rightarrow\lb 2=1$ bit，正好等于“一次公平抛硬币”的信息量；
比值 $=4\Rightarrow\lb 4=2$ bit；比值 $=\frac14\Rightarrow\lb\frac14=-2$ bit。
所以你读这个数可以这样读：\textbf{“$1$ bit $=$ 一次公平抛硬币的把握”}。

\subsubsection{第一个数值例子：天气与带伞}

\begin{calc}{例 1：把天气表的四格全算出来}
学长给了一张真实点的 $2\times2$ 表（横向是带不带伞，纵向是下不下雨，
格子里的数是概率）。先看\textbf{边缘}（每行、每列加起来）：

\begin{center}
\small
\begin{tabular}{@{}lccc@{}}
\toprule
 & 带伞 $y=u$ & 不带伞 $y=\bar u$ & 边缘 $p(x)$ \\
\midrule
下雨 $x=r$ & $0.4$ & $0.1$ & $0.5$ \\
不下雨 $x=\bar r$ & $0.2$ & $0.3$ & $0.5$ \\
边缘 $p(y)$ & $0.6$ & $0.4$ & $1$ \\
\bottomrule
\end{tabular}
\end{center}

\noindent\textbf{四格逐个算}（分子是联合概率，分母是两个边缘概率之积）：
\[
\begin{aligned}
\text{第 1 格：}\ \I{r;u}&=\lb\frac{p(r,u)}{p(r)p(u)}=\lb\frac{0.4}{0.5\times0.6}
=\lb\frac{0.4}{0.3}=\lb\frac{4}{3}=\lb 4-\lb 3\\
&=2-1.58496=0.41504\approx0.415\bit ,\\
\text{第 2 格：}\ \I{r;\bar u}&=\lb\frac{p(r,\bar u)}{p(r)p(\bar u)}
=\lb\frac{0.1}{0.5\times0.4}=\lb\frac{0.1}{0.2}=\lb\frac12=-1\bit ,\\
\text{第 3 格：}\ \I{\bar r;u}&=\lb\frac{p(\bar r,u)}{p(\bar r)p(u)}
=\lb\frac{0.2}{0.5\times0.6}=\lb\frac{0.2}{0.3}=\lb\frac23\\
&=\lb 2-\lb 3=1-1.58496=-0.58496\approx-0.585\bit ,\\
\text{第 4 格：}\ \I{\bar r;\bar u}&=\lb\frac{p(\bar r,\bar u)}{p(\bar r)p(\bar u)}
=\lb\frac{0.3}{0.5\times0.4}=\lb\frac{0.3}{0.2}=\lb\frac32\\
&=\lb 3-\lb 2=1.58496-1=0.58496\approx0.585\bit .
\end{aligned}
\]
怎么读这四个数：
\begin{itemize}[leftmargin=1.4em,itemsep=1pt,topsep=2pt]
  \item 第 1 格 $+0.415$：下雨时带伞的机会是“不管天气瞎猜”的 $\frac43$ 倍——
        大爷那句“下雨”确实把你往“带伞”这边推了一把。
  \item 第 2 格 $-1$：\textbf{负数！}“下雨”这条消息让“不带伞”\textbf{更难}发生了
        （$0.4$ 降到 $0.2$），所以是 $-1$ bit 的“反向帮助”。负号不是算错。
  \item 第 3、4 格是一对反号：不下雨时“带伞”变难（$-0.585$）、
        “不带伞”变容易（$+0.585$），符合常识。
\end{itemize}
小结：$+0.415,\,-1,\,-0.585,\,+0.585$——有正有负，说明“带伞”这件事
有时帮我们猜天气、有时反而误导（不带伞的人里也有相当多是懒，
跟天气没关系）。
\end{calc}

\begin{calc}{例 2：掷两颗骰子——完全独立，答案就是 $0$}
$X$ 是第一颗骰子的点数，$Y$ 是第二颗。取 $x=6,y=6$：
\[
p(x,y)=\frac1{36},\quad p(x)=\frac16,\quad p(y)=\frac16
\ \Longrightarrow\
\I{x;y}=\lb\frac{1/36}{(1/6)\times(1/6)},
\]
\[
\Longrightarrow\ \I{x;y}=\lb\frac{1/36}{1/36}=\lb 1=0\bit .
\]
两颗骰子各管各的，“第二颗是几点”对“第一颗是几点”一分钱信息也不给，
这就是 $0$ 的含义。
\end{calc}

\begin{calc}{例 3：掷两颗骰子——有牵连时是正数}
练一个“看着没关系、其实强关联”的：设 $X$ 是\textbf{第一颗骰子的点数}
（取值 $1,\dots,6$），$Y$ 是\textbf{两颗骰子的点数之和}（取值 $2,\dots,12$）。
现在看两个事件：$x$ 表示“第一颗是 $6$ 点”，$y$ 表示“和是 $12$”。
注意“和是 $12$”只有一种可能：$6+6$。于是
\[
p(x=6)=\frac16,\qquad p(y=12)=\frac1{36},\qquad p(x=6,\,y=12)=\frac1{36},
\]
\[
\I{x=6;\,y=12}=\lb\frac{1/36}{(1/6)\times(1/36)}
=\lb\frac{1/36}{1/216}=\lb 6\approx2.585\bit .
\]
人话：一旦知道“和是 $12$”，第一颗必然是 $6$——这条消息几乎把答案
直接端到你手里，所以是很大的正数 $2.585$ bit。
顺带看一眼上界：事件 $x=6$ 自己的信息量是 $I(x=6)=-\lb\frac16=\lb 6\approx2.585$ bit，
这里 $\I{x;y}=2.585$ 正好顶到 $I(x)$，说明“全知道了”。
（与此对比，例 1 里概率高达 $0.4$ 的格子只值 $0.415$ bit，
概率只有 $0.0278$ 的这一格却值 $2.585$ bit——\textbf{概率大小和关联强度
根本不是一回事}，这一点问答部分还会说。）
\end{calc}

\subsubsection{和自信息的关系：三种等价的写法}

把 $p(x,y)=p(x|y)p(y)$ 代进定义，再把对数拆开（$\lb$ 把乘积变加法，
这又是本讲那个“技巧”），立刻得到
\[
\I{x;y}=\lb\frac{p(x|y)}{p(x)}
=\underbrace{-\lb p(x)}_{I(x)}-\underbrace{\bigl(-\lb p(x|y)\bigr)}_{I(x|y)}
=I(x)-I(x|y),
\]
其中 $p(x|y)$ 读作“\textbf{在 $y$ 已经发生的条件下} $x$ 发生的概率”，
$I(x|y)=-\lb p(x|y)$ 叫\textbf{条件自信息}（知道 $y$ 以后 $x$ 还剩多少信息量）。
同理
\[
\I{x;y}=I(y)-I(y|x)=I(x)+I(y)-I(x,y).
\]
人话翻译：\textbf{互信息 $=$ 原来需要的“惊讶度”，减去知道 $y$ 之后剩下的“惊讶度”。}
这已经是下一节“平均互信息”三条公式的雏形了——它们是同一件事
在“事件层面”和“随机变量层面”的两次出现。

% ==================================================================
\subsection{互信息可正、可负、可为零；以及它最可爱的性质：对称}

\begin{defbox}{符号法则（一定要记住这一条）}
\begin{center}
\small
\begin{tabular}{@{}lll@{}}
\toprule
情形 & 数学条件 & 人话 \\
\midrule
正 & $p(x|y)>p(x)$ & 知道 $y$ 之后，$x$ 更可能发生了（帮上忙） \\
零 & $p(x|y)=p(x)$，即 $p(x,y)=p(x)p(y)$ & 知不知道 $y$，对猜 $x$ 毫无影响 \\
负 & $p(x|y)<p(x)$ & 知道 $y$ 之后，$x$ 反而更不可能了（帮倒忙） \\
\bottomrule
\end{tabular}
\end{center}
注意“零”那一行的\textbf{充要}关系：对这一对事件而言，
$\I{x;y}=0\iff x$ 与 $y$ \textbf{独立}。为什么？因为 $\lb$ 只在自变量为 $1$ 时取 $0$。
\end{defbox}

\begin{idea}{为什么会“帮倒忙”（负数不是bug）}
例 1 里的 $\I{r;\bar u}=-1$ bit 可以这么理解：你已经知道“今天下雨”了，
按常识下雨更可能带伞，那么“不带伞”的可能性就从 $0.4$ 掉到 $0.2$。
一条消息让某件事\emph{更不可能}，它就在这件事上提供了负的信息量
——就像你听到“今天有暴雨”，会更确定“今天不适合晾被子”。

而“平均互信息”永远非负（后面“基本性质”一节证明）：\textbf{负的格子会被正的格子抵消}
（例 1 里 $-1$ 被 $+0.415$、$+0.585$ 抵掉了）。
所以考试里如果算出 $\I{X;Y}<0$，一定是算错了。
\end{idea}

\begin{defbox}{对称性}
\[
\I{x;y}=\lb\frac{p(x,y)}{p(x)p(y)}
=\lb\frac{p(y,x)}{p(y)p(x)}=\I{y;x}.
\]
人话：“知道下雨对猜带伞的帮助”和“知道带伞对猜下雨的帮助”是\textbf{同一个数}。
\end{defbox}

\begin{idea}{对称性为什么值得单独说？因为它最容易在考试里被写反}
互信息对称，但\textbf{熵不对称}：$H(X|Y)\ne H(Y|X)$，一般是不等的！
例 1 里 $H(X|Y)=0.87549$ 而 $H(Y|X)=0.84644$，两者差 $0.029$。
记住这一对“长得像但不是一回事”的量：
\begin{itemize}[leftmargin=1.4em,itemsep=1pt,topsep=1pt]
  \item $\I{X;Y}$：两个圆\textbf{重叠}的面积，对称；
  \item $H(X|Y)$：$X$ 那个圆里\textbf{没被 $Y$ 盖住}的面积，不对称。
\end{itemize}
\end{idea}

还有两个上界，考试常拿来做“不可能性”判断：因为 $p(x|y)\le1$，所以
\[
\I{x;y}=I(x)-I(x|y)\le I(x),\qquad
\I{x;y}\le I(y).
\]
人话：一件事带给你的关于 $x$ 的信息，不可能超过“$x$ 自己发生”所含的信息量
（例 3 里 $\I{x;y}=2.585=I(x)$ 就是“全知道”的临界点）。
这两条上界，就是后面“基本性质”一节里极值性 $\I{X;Y}\le\min\{\Hx{X},\Hx{Y}\}$ 的事件级原型。

% ==================================================================
\subsection{平均互信息：把每一格的账按“发生频率”加权}

事件互信息只管\textbf{一个格子}。但我们通常要问的是整体：
“变量 $X$ 和变量 $Y$ 到底有多相关？”于是把所有格子的账按
\textbf{它实际发生的频率}加权平均。这个“加权平均”就是数学期望 $\E[\cdot]$。

\begin{defbox}{定义（平均互信息，教材定义 3.1.3）}
\[
\I{X;Y}=\E\bigl[\I{x;y}\bigr]
=\sum_{x\in\mathcal X}\sum_{y\in\mathcal Y} p(x,y)\,\lb\frac{p(x,y)}{p(x)p(y)} .
\]
人话：\textbf{对每一格算一句“帮了多少”，再按那一格实际出现的概率
$p(x,y)$ 加权平均}。权重是 $p(x,y)$ 而不是 $p(x)p(y)$——
因为我们算的是“现实中真的发生过的那些情况”的平均感受。
\end{defbox}

\begin{qa}{为什么权重是 $p(x,y)$，不是 $p(x)p(y)$？}
因为期望的定义是“取值 $\times$ 该取值出现的概率”。$\I{x;y}$ 这个量
出现在“$x,y$ 同时发生”这个场合，而“$x,y$ 同时发生”的概率正是 $p(x,y)$。
用 $p(x)p(y)$ 当权重是在算一个\textbf{虚构}的平均：它按“假设两件事已经独立”
来分配频次，而不是按现实里的频次，于是“$X$ 常见还是少见”这个信息会被一起抹掉，
算出来就不再是“知道 $Y$ 到底帮我多少”了。
一句话：\textbf{对数里面已经用 $p(x,y)$ 算账，外面的权重也得用同一个频率，
内外一致才叫平均互信息}。
\end{qa}

还有一种“信道视角”的写法，把 $p(x,y)=p(x)p(y|x)$ 代进去：
\[
\I{X;Y}=\sum_{x}p(x)\sum_{y}p(y|x)\,\lb\frac{p(y|x)}{p(y)}
=\sum_{x}p(x)\sum_{y}p(y|x)\lb\frac{p(y|x)}{\sum_{x'}p(x')p(y|x')},
\]
其中 $p(y|x)$ 读作“\textbf{输入是 $x$ 时，输出变成 $y$ 的概率}”（后面章节的
“转移概率矩阵”就是它）。课本（式 3.1.16）之所以要写这一版，是因为它说明了：
\textbf{平均互信息完全由“输入分布 $p(x)$ $+$ 信道 $p(y|x)$”决定}——
这就是第 4、5 章信道容量讨论的出发点。

% ==================================================================
\subsection{平均互信息的三种等价写法（三条都要背，考试任选一条）}

\begin{defbox}{定理（互信息与其他熵的关系，教材定理 3.2.1、3.2.2、3.2.3）}
\[
\I{X;Y}=\Hx{X}-\Hx{X|Y}=\Hx{Y}-\Hx{Y|X}=\Hx{X}+\Hx{Y}-\Hx{X,Y}.
\]
\end{defbox}

三条公式的人话翻译（这段是本讲最值钱的部分，请逐字读）：

\begin{itemize}[leftmargin=1.6em,itemsep=3pt,topsep=2pt]
  \item \textbf{第一条 $\I{X;Y}=\Hx{X}-\Hx{X|Y}$：}
        \textbf{互信息 $=$ 原来的不确定性 $-$ 知道 $Y$ 之后剩下的不确定性。}
        这是最本色的一条，就是本讲开头那句话的公式版：
        “不确定性减少了多少”。也正因为 $\Hx{X|Y}\ge0$（剩下的不确定性
        不可能是负的），立刻得到 $\I{X;Y}\le\Hx{X}$。
  \item \textbf{第二条 $\I{X;Y}=\Hx{Y}-\Hx{Y|X}$：}
        \textbf{反过来问也一样}：知道 $X$ 之后，$Y$ 的不确定性减少量
        和上一条\textbf{完全相等}。这就是平均值层面的对称性
        （事件级对称 $\I{x;y}=\I{y;x}$ 加权平均后自然还对称）。
  \item \textbf{第三条 $\I{X;Y}=\Hx{X}+\Hx{Y}-\Hx{X,Y}$：}
        \textbf{两个各自的不确定性加起来，减去“合起来看”的不确定性。}
        为什么？如果 $X,Y$ 完全没关系，合起来看的不确定性正好是两者相加
        （$H(X,Y)=H(X)+H(Y)$），多出来的部分就是 $0$；两者越相关，
        “合起来看”越是省事（重复的信息只算一次），$H(X,Y)$ 就越小，
        差额 $\I{X;Y}$ 就越大。\textbf{所以 $\I{X;Y}$ 就是“重复计算的那部分不确定性”。}
\end{itemize}

\begin{idea}{韦恩图记忆法（教材图 3.2 就是它）}
画两个圆：左边圆面积 $\Hx{X}$，右边圆面积 $\Hx{Y}$，两圆并起来
$\Hx{X,Y}$，中间重叠部分是 $\I{X;Y}$。于是
\begin{itemize}[leftmargin=1.4em,itemsep=1pt,topsep=1pt]
  \item 左边“月牙” $=$ $\Hx{X|Y}$：$X$ 独有的不确定性（$Y$ 没帮你解决的那部分）；
  \item 右边“月牙” $=$ $\Hx{Y|X}$；
  \item 重叠 $=$ $\I{X;Y}$：双方共有的、重复的那部分不确定性；
  \item 并集 $=$ $\Hx{X,Y}$；
  \item 两圆不相交 $\iff$ 独立 $\iff$ $\I{X;Y}=0$。
\end{itemize}
考试时在草稿纸角上画这两个圈，$7$ 条形如“$\Hx{X|Y}=\Hx{X,Y}-\Hx{Y}$”
的关系式全都自动想得起来。
\end{idea}

\begin{calc}{顺手验证“等价”不是口号：用定义推第一条}
学长不想让你背，想让你看它为什么成立。从定义
$\I{X;Y}=\sum_{x,y}p(x,y)\lb\frac{p(x,y)}{p(x)p(y)}$ 出发，
把对数里的分式拆成 $p(x|y)/p(x)$（\textbf{技巧：}因为 $p(x,y)=p(x|y)p(y)$，
分式上下同除 $p(y)$ 即可）：
\[
\I{X;Y}=\sum_{x,y}p(x,y)\lb p(x|y)-\sum_{x,y}p(x,y)\lb p(x).
\]
第一项 $=\sum_{y}p(y)\sum_{x}p(x|y)\lb p(x|y)=-\Hx{X|Y}$；
第二项 $=\sum_{x}\Bigl(\sum_y p(x,y)\Bigr)\lb p(x)=\sum_x p(x)\lb p(x)=-\Hx{X}$。
两式相减即 $\I{X;Y}=\Hx{X}-\Hx{X|Y}$。$\square$
（看到没——\textbf{每一次“把互信息换成一堆 $H$ 的加减”都是这个套路}：
先把 $\lb$ 里的分式拆成“联合 $\div$ 边缘”，再让一组求和“收”成边缘分布。）
\end{calc}

% ==================================================================
\subsection{手把手：把天气表从头算到尾，三条公式对答案}

\begin{calc}{完整计算（$2\times2$ 表，算术不跳步）}
\textbf{已知}（就是例 1 那张表）：
\[
p(r,u)=0.4,\quad p(r,\bar u)=0.1,\quad p(\bar r,u)=0.2,\quad p(\bar r,\bar u)=0.3 .
\]
换算器（就是本讲开头要你背的那三个数）：$\lb 3=1.58496$，$\lb 5=2.32193$，$\lb 10=3.32193$。
由它们立马得到：$\lb 0.4=-1.32193$，$\lb 0.6=-0.73697$，
$\lb 0.1=-3.32193$，$\lb 0.2=-2.32193$，$\lb 0.3=-1.73697$。

\medskip
\noindent\textbf{步骤 1：先算边缘分布}（行、列分别相加）
\[
p(r)=0.4+0.1=0.5,\quad p(\bar r)=0.2+0.3=0.5,\quad
p(u)=0.4+0.2=0.6,\quad p(\bar u)=0.1+0.3=0.4 .
\]
$X$ 是公平的（一半一半），$Y$ 偏一点（$60\%$ 带伞）。

\medskip
\noindent\textbf{步骤 2：算三个熵}（都用 $H=-\sum p\lb p$）
\[
\begin{aligned}
\Hx{X}&=-0.5\lb 0.5-0.5\lb 0.5=0.5\times1+0.5\times1=1.00000\bit ,\\
\Hx{Y}&=-0.6\lb 0.6-0.4\lb 0.4
=0.6\times0.73697+0.4\times1.32193=0.97095\bit ,\\
\Hx{X,Y}&=-\bigl[\,0.4\lb 0.4+0.1\lb 0.1+0.2\lb 0.2+0.3\lb 0.3\,\bigr]\\
&=0.4\times1.32193+0.1\times3.32193+0.2\times2.32193+0.3\times1.73697\\
&=0.52877+0.33219+0.46439+0.52109=1.84644\bit .
\end{aligned}
\]

\medskip
\noindent\textbf{步骤 3：用三条公式分别算 $\I{X;Y}$，看结果一不一样}
\[
\begin{aligned}
\text{式一：}\ &\Hx{X|Y}=\Hx{X,Y}-\Hx{Y}=1.84644-0.97095=0.87549,\\
&\I{X;Y}=\Hx{X}-\Hx{X|Y}=1.00000-0.87549=0.12451\bit ;\\
\text{式二：}\ &\Hx{Y|X}=\Hx{X,Y}-\Hx{X}=1.84644-1.00000=0.84644,\\
&\I{X;Y}=\Hx{Y}-\Hx{Y|X}=0.97095-0.84644=0.12451\bit ;\\
\text{式三：}\ &\I{X;Y}=\Hx{X}+\Hx{Y}-\Hx{X,Y}
=1.00000+0.97095-1.84644=0.12451\bit .
\end{aligned}
\]
三条全给 $0.12451$，一bit不差。

\medskip
\noindent\textbf{步骤 4：用定义（四格加权）再算一遍，第四次对答案}
\[
\begin{aligned}
\I{X;Y}={}&0.4\times\lb\frac{0.4}{0.3}+0.1\times\lb\frac{0.1}{0.2}
+0.2\times\lb\frac{0.2}{0.3}+0.3\times\lb\frac{0.3}{0.2}\\
={}&0.4\times0.41504+0.1\times(-1)+0.2\times(-0.58496)+0.3\times0.58496\\
={}&0.16602-0.10000-0.11699+0.17549=0.12452 .
\end{aligned}
\]
（最后一格的 $0.00001$ 差是四舍五入攒的，不是公式错。
考试里一般保留 $3\sim4$ 位有效数字即可。）

\medskip
\noindent\textbf{步骤 5：这个 $0.1245$ bit 说明什么？}
$\Hx{X}=1$ bit 是“完全猜不着”，$\Hx{X|Y}=0.8755$ bit 只降了
$0.1245$ bit（约 $12\%$）。人话：\textbf{光知道“我带不带伞”，
对你猜“下不下雨”的帮助很有限}——因为带伞这事还掺了“懒”
“教室有伞”等与天气无关的因素。对应地，
$\I{X;Y}=0.12451<0.97095=\Hx{Y}<1=\Hx{X}$，两个上界都满足。

\noindent\textbf{顺手做个“极值”对照}（这就是下一小节要说的性质）：
\begin{itemize}[leftmargin=1.4em,itemsep=1pt,topsep=1pt]
  \item \textbf{完全同步的表} $p(r,u)=0.6,\ p(\bar r,\bar u)=0.4$，另两格为 $0$
        （下雨一定带伞、不下雨一定不带）：两边都是 $\Hx{0.6}=0.97095$ bit，
        $\Hx{X|Y}=0$，故 $\I{X;Y}=0.97095\bit$，正好顶到上界。
  \item \textbf{独立表} $p(r,u)=0.3,\ p(r,\bar u)=0.2,\ p(\bar r,u)=0.3,\
        p(\bar r,\bar u)=0.2$：每格都等于两边缘之积，故 $\I{X;Y}=\lb1=0\bit$。
\end{itemize}
\textbf{注意独立表和本节的表边缘分布完全一样}（$X$ 都 $0.5/0.5$，$Y$ 都 $0.6/0.4$），
互信息却一个是 $0$、一个是 $0.12451$——\textbf{光看边缘定不出互信息，
关键看“怎么搭配”}。
\end{calc}

\begin{calc}{第二组数据：二元对称信道（BSC）——考试最爱的小例子}
$X$ 是一枚公平硬币（$0/1$ 各 $\frac12$）。$Y$ 是对 $X$ 的观测，
但每次有 $0.1$ 的概率看错（专业名称：错误转移概率 $p=0.1$）。
于是 $p(y|x)$ 两行都是 $(0.9,\,0.1)$ 的两个排列，且
\[
p(y=0)=0.5\times0.9+0.5\times0.1=0.5,
\]
即 $Y$ 也是公平硬币，$\Hx{Y}=1$ bit；而（两行条件分布相同，直接加权）
\[
\Hx{Y|X}=0.5\bigl[-\!0.9\lb0.9-0.1\lb0.1\bigr]\times2
=-0.9\lb 0.9-0.1\lb 0.1 .
\]
用换算器：$\lb 0.9=\lb\frac{9}{10}=\lb 9-\lb 10$，而 $\lb 9=2\lb 3=3.16993$，
故 $\lb 0.9=3.16993-3.32193=-0.15200$；又 $\lb 0.1=-3.32193$。所以
\[
\Hx{Y|X}=0.9\times0.15200+0.1\times3.32193=0.13680+0.33219=0.46899\bit ,
\]
\[
\I{X;Y}=\Hx{Y}-\Hx{Y|X}=1-0.46899=0.53101\bit .
\]
人话：$1$ bit 的不确定性，信道只送到 $0.531$ bit，剩下 $0.469$ bit
就是“噪声吃掉的部分”。这条式子后面第 5 章讲信道容量时天天用。
（顺便记住 $\Hx{0.1}\approx0.469$ 这个数，考试出现频率极高。）
\end{calc}

% ==================================================================
\subsection{互信息的四条基本性质（每条一句人话）}

\begin{defbox}{性质清单}
设 $X,Y$ 为离散随机变量，则
\begin{enumerate}[leftmargin=1.8em,itemsep=2pt,topsep=2pt]
  \item \textbf{非负性：}$\I{X;Y}\ge0$，且等号成立 $\iff$ $X,Y$ 统计独立。
        \emph{人话：知道点东西总比什么都不知道好；平均值上不可能“帮倒忙”。}
  \item \textbf{极值性：}$\I{X;Y}\le\min\{\Hx{X},\Hx{Y}\}$。
        \emph{人话：从 $Y$ 里能榨出的关于 $X$ 的信息，不可能超过 $X$ 本身的不确定性，
        也不可能超过“$Y$ 这个容器里总共装的东西”。}
  \item \textbf{独立为零：}$X,Y$ 独立 $\iff\I{X;Y}=0$（与性质 1 的等号条件同一件事）。
        \emph{人话：两件互不相干的事，互相提供不了任何信息。}
  \item \textbf{自信息：}$\I{X;X}=\Hx{X}$。
        \emph{人话：自己当然完全知道自己，不确定性一次全消。}
\end{enumerate}
\end{defbox}

\subsubsection{非负性的证明（只用 $\ln u\le u-1$ 这一个技巧）}

\begin{calc}{证明 $\I{X;Y}\ge0$，并拿到取等条件}
约定：$p(x,y)=0$ 的格子那项算 $0$（因为 $\lim_{t\to0}t\lb t=0$），
所以只需处理 $p(x,y)>0$ 的格子。\textbf{技巧：用一个初等不等式}
\[
\ln u\le u-1\quad(u>0)\ \Longleftrightarrow\ -\lb u\ \ge\ \frac{1-u}{\ln 2}
\quad(\text{因为 } \lb u=\tfrac{\ln u}{\ln 2},\ \ln 2\approx0.6931).
\]
把 $u=\dfrac{p(x)p(y)}{p(x,y)}$ 代进去：
\[
-\lb\frac{p(x)p(y)}{p(x,y)}\ \ge\ \frac{1}{\ln 2}\left(1-\frac{p(x)p(y)}{p(x,y)}\right).
\]
两边乘 $p(x,y)$（正数，不等号方向不变）并对所有格子求和：
\[
\I{X;Y}=\sum_{x,y}p(x,y)\lb\frac{p(x,y)}{p(x)p(y)}
\ \ge\ \frac{1}{\ln 2}\left[\sum_{x,y}p(x,y)-\sum_{x,y}p(x)p(y)\right].
\]
括号里两个和都是 $1$（分别是“全表加起来 $=1$”和“边缘乘起来分块求和 $=1$”），
所以右边 $=\dfrac{1}{\ln 2}(1-1)=0$，即 $\I{X;Y}\ge0$。

\textbf{取等条件}：$\ln u=u-1$ 只在 $u=1$ 时成立，即每个非零格子都满足
$p(x,y)=p(x)p(y)$——这正是\textbf{独立}。反过来独立时每格比值都是 $1$，
每项都是 $\lb1=0$。证完。$\square$
\end{calc}

\subsubsection{极值性与另外两条}

\begin{itemize}[leftmargin=1.6em,itemsep=3pt,topsep=2pt]
  \item \textbf{$\I{X;Y}\le\Hx{X}$}：由 $\I{X;Y}=\Hx{X}-\Hx{X|Y}$，
        而条件熵 $\Hx{X|Y}\ge0$（\emph{它是一堆“$-\lb p$”的加权平均，每项都非负}），
        所以自动成立；取等 $\iff\Hx{X|Y}=0$，即知道 $Y$ 就等于知道 $X$。
  \item \textbf{$\I{X;Y}\le\Hx{Y}$}：把上一条里的 $X,Y$ 换个名字（因为对称）。
        两条合起来就是\textbf{极值性}：
        \[
        \I{X;Y}\le\min\{\Hx{X},\Hx{Y}\}.
        \]
  \item \textbf{$\I{X;X}=\Hx{X}-\Hx{X|X}=\Hx{X}-0=\Hx{X}$}：
        用第一条件公式：已知 $X$ 时“$X$ 还是什么”已经确定了，条件概率
        $p(x|x)=1$，故 $-\lb1=0$，也就是“剩下的不确定性 $=0$”。数值感受：BSC 例子里若不出错
        （$p=0$），则 $\I{X;Y}=\Hx{Y}=\Hx{X}=1$ bit，与 $\I{X;X}$ 一样大。
\end{itemize}

\begin{exam}{两条性质的高频用法}
\begin{enumerate}[leftmargin=1.8em,itemsep=2pt,topsep=2pt]
  \item 题目问“$X,Y$ 是否独立”$\Rightarrow$ 只看一个数：$\I{X;Y}$ 是否为 $0$
        （等价地 $p(x,y)$ 是否处处 $=p(x)p(y)$）。
  \item 题目要你证明某个量“$\le\Hx{X}$”$\Rightarrow$ 把那个量写成
        $\I{X;Y}$ 或 $\Hx{X}-\Hx{X|Y}$，再用条件熵非负。
  \item 让你“证明 $\I{X;Y}\ge0$”就是上面那段不等式，或者干脆用
        $\I{X;Y}=\Hx{X}-\Hx{X|Y}\ge0$（两种写法阅卷都认）。
\end{enumerate}
\end{exam}

% ==================================================================
\subsection{条件互信息与链式法则：把“知道”分成两步}

现实里的“知道”常常是一步一步来的：先知道 $Y$，再知道 $Z$。
条件互信息问的是：\textbf{在已经知道 $Y$ 的前提下，$Z$ 还能额外告诉你多少关于 $X$ 的事？}

\begin{defbox}{定义（条件互信息，教材式 3.3.3）}
\[
\I{X;Y|Z}=\sum_{z}p(z)\sum_{x,y}p(x,y|z)\,\lb\frac{p(x,y|z)}{p(x|z)\,p(y|z)} .
\]
人话：\textbf{把 $Z$ 固定成某一个值 $z$，在这一层里算一遍普通的互信息，
再按 $p(z)$ 加权平均}。所以条件互信息不是新东西，就是“分组算互信息，
然后组间加权平均”。三条等价写法照抄上一节（把每个 $H$ 都加上“$|Z$”）：
\[
\I{X;Y|Z}=\Hx{X|Z}-\Hx{X|YZ}=\Hx{Y|Z}-\Hx{Y|XZ}
=\Hx{X|Z}+\Hx{Y|Z}-\Hx{X,Y|Z}.
\]
同样有非负性 $\I{X;Y|Z}\ge0$，和对称性 $\I{X;Y|Z}=\I{Y;X|Z}$。
\end{defbox}

\begin{idea}{“分两步知道”的直觉}
$X$ 是一台机器的故障原因，$Y$ 是仪表读数，$Z$ 是维修手册。
\begin{itemize}[leftmargin=1.4em,itemsep=1pt,topsep=1pt]
  \item $\I{X;Y}$：光看仪表，能推断多少故障原因；
  \item $\I{X;Z|Y}$：\textbf{已经看过仪表之后}再翻手册，还能额外推断出多少；
  \item $\I{X;YZ}$：仪表 $+$ 手册一起用，总共能推断多少。
\end{itemize}
“一共的收获”$=$“仪表给的”$+$“看过仪表之后再翻手册额外给的”——
这句话本身就是加法的定义，不用证明。
把这句大白话写成公式，就是链式法则（下面这个式子）。
\end{idea}

\begin{defbox}{链式法则（教材式 3.3.9，作业 3.3 考的就是它）}
\[
\I{X;YZ}=\I{X;Y}+\I{X;Z|Y}\ =\ \I{X;Z}+\I{X;Y|Z}.
\]
人话：\textbf{（$Y$, $Z$ 一起提供的关于 $X$ 的信息）
$=$（$Y$ 提供的）$+$（已知 $Y$ 后 $Z$ 额外提供的）}。
两个式子都对，因为“先知道 $Y$ 还是先知道 $Z$”只是顺序不同，总量一样。
\end{defbox}

\begin{calc}{证明链式法则（只有三行，但每行都要会对）}
记 $\I{X;YZ}$ 就是“把 $(Y,Z)$ 当成一个整体”再用公式：
\[
\I{X;YZ}=\Hx{X}-\Hx{X|YZ}
=\underbrace{\bigl[\Hx{X}-\Hx{X|Y}\bigr]}_{\I{X;Y}}
+\underbrace{\bigl[\Hx{X|Y}-\Hx{X|YZ}\bigr]}_{\I{X;Z|Y}} .
\]
第一个方括号按定义就是 $\I{X;Y}$；第二个方括号正是“已知 $Y$ 之后，
再把 $Z$ 弄明白所减少的不确定性”，也就是 $\I{X;Z|Y}$（见上面条件互信息
的第一条等价写法）。$\square$

\textbf{注意那个“技巧”}：这里没有算任何概率，只是加一项减一项
（$-\Hx{X|Y}$ 和 $+\Hx{X|Y}$），然后按定义认领。证明熵关系的题都是这一招：
\textbf{加一项减一项，凑成定义的样子}。
\end{calc}

\begin{calc}{两个数值例子（一个“白给”，一个“惊人”）}
\textbf{例 A（冗余：第二份只是重复）}：$X$ 是公平硬币，$Y=X$，$Z=X$（两份副本）。
$\I{X;Y}=\Hx{X}=1$ bit；但知道 $Y$ 之后 $Z$ 一点新东西都没有，
所以 $\I{X;Z|Y}=0$；链式法则给出 $\I{X;YZ}=1+0=1$ bit。
人话：\textbf{同一个消息听两遍，第二遍不值钱。}

\textbf{例 B（协同：$1=0+1$，这就是作业 3.1 的标准答案）}：
$X,Y$ 是两枚\textbf{独立}的公平硬币，令 $Z=X\oplus Y$（异或：相同为 $0$，不同为 $1$）。
\begin{itemize}[leftmargin=1.4em,itemsep=1pt,topsep=1pt]
  \item 先算 $\I{X;Y}$：因为独立，$p(x,y)=p(x)p(y)$ 处处成立，所以 $\I{X;Y}=0$ bit。
  \item 再算 $\I{X;Y|Z}$：固定 $z=0$，此时 $(x,y)$ 只能是 $(0,0)$ 或 $(1,1)$，
        各占一半；$p(x|z=0)=\tfrac12$，$p(y|z=0)=\tfrac12$，
        $p(x,y|z=0)=\tfrac12$。于是这一组里
        $\I{x;y|z=0}=\lb\frac{1/2}{(1/2)(1/2)}=\lb2=1$ bit。
        固定 $z=1$ 同理也是 $1$ bit。加权平均：$\I{X;Y|Z}=1$ bit。
  \item 验证链式法则：$\I{X;YZ}=\I{X;Y}+\I{X;Z|Y}=0+\I{X;Z|Y}$。
        而由于 $X=Y\oplus Z$，知道 $(Y,Z)$ 就等于知道 $X$，
        所以 $\I{X;YZ}=\Hx{X}=1$ bit，从而 $\I{X;Z|Y}=1$ bit。
        于是 $1=0+1$，两边相等。
\end{itemize}
人话：\textbf{两枚硬币本来毫无关系（互信息 $0$），但一旦告诉你“它俩的异或”，
它们就锁死了（互信息 $1$）}。这说明“$X$ 与 $Y$ 独立”和
“已知 $Z$ 时 $X$ 与 $Y$ 独立”是两回事——这是下一节那个反例的伏笔。
\end{calc}

% ==================================================================
\subsection{数据处理：信息处理只会丢信息，不会凭空造信息}

\subsubsection{马尔可夫链是什么意思}

$X\to Y\to Z$ 构成\textbf{马尔可夫链}，意思是：
\textbf{$Z$ 只通过 $Y$ 和过去相连}，用条件概率写就是
\[
p(z|x,y)=p(z|y)\quad\text{对所有 }x,y,z ,
\]
读作：“已经知道 $Y$ 之后，$X$ 的取值对 $Z$ 没有任何额外影响”。
生活版：$X$ 是你写的作业，$Y$ 是拍照后的照片，$Z$ 是学长的转述。
照片坏了以后，学长只能照着照片说——\textbf{他说的只取决于照片，
不取决于你原件写了什么}（原件他没见过）。

\begin{defbox}{数据处理定理（教材习题 3.10，PPT 讲评重点）}
若 $X\to Y\to Z$ 是马尔可夫链，则
\[
\I{X;Z}\le\I{X;Y}.
\]
人话：\textbf{信息处理（拍照、压缩、转发、算平均）只会让“关于 $X$ 的信息”
变少或不变，绝不可能变多}——处理不创造信息，只损失信息。
\end{defbox}

\begin{calc}{三步证明（这题考试的写法就这样）}
\textbf{第 1 步}（链式法则，两种顺序都写出来）：
\[
\I{X;YZ}=\I{X;Y}+\I{X;Z|Y}=\I{X;Z}+\I{X;Y|Z}.
\]
\textbf{第 2 步}（用马尔可夫条件，把一项干掉）：
因为 $p(z|x,y)=p(z|y)$，所以“已知 $Y$ 时，$X$ 与 $Z$ 独立”，
从而 $\I{X;Z|Y}=0$。于是第 1 步的左边两个式子化为
\[
\I{X;Y}=\I{X;YZ}=\I{X;Z}+\I{X;Y|Z}.
\]
\textbf{第 3 步}（用条件互信息非负性收尾）：
$\I{X;Y|Z}\ge0$，故 $\I{X;Y}\ge\I{X;Z}$。$\square$

\textbf{记忆口诀}：\textbf{“链式法则 $+$ 马尔可夫杀一项 $+$ 非负性收尾”}，
三道工序，全章证明题的通用模板。
\end{calc}

\begin{calc}{数值感受：两次拍照，信息一路变少（BSC 级联）}
接着用前面“二元对称信道”那个例子：$X$ 公平硬币 $\to$ 拍照（错 $0.1$）得 $Y$ $\to$ 再拍一次
（错 $0.1$）得 $Z$。
\begin{itemize}[leftmargin=1.4em,itemsep=1pt,topsep=1pt]
  \item 第一步：$\I{X;Y}=1-0.46899=0.53101$ bit（前面算过）。
  \item 两次拍照合起来等价于“一次错多少”？做一次“翻错几次”的分类：
        只看错一次（第一次错且第二次对，或第一次对且第二次错）
        $=0.1\times0.9+0.9\times0.1=0.18$；两次都对 $=0.9\times0.9=0.81$；
        两次都错 $=0.1\times0.1=0.01$。从头到尾看错一次的概率就是 $0.18$，
        所以 $X\to Z$ 等价于一个错误率 $0.18$ 的 BSC。
  \item 于是 $\Hx{Z|X}=-[0.18\lb0.18+0.82\lb0.82]$。拆 $\lb$ 的\textbf{技巧：
        把小数写成“分母含 $50$ 的分数”}：
        $\lb0.18=\lb\frac{9}{50}=\lb9-\lb50=3.16993-5.64386=-2.47393$
        （其中 $\lb 9=2\lb 3$，$\lb 50=\lb2+2\lb5$），
        同理 $\lb0.82=\lb\frac{41}{50}=5.35755-5.64386=-0.28631$。代入：
        \[
        \Hx{Z|X}=0.18\times2.47393+0.82\times0.28631
        =0.44531+0.23477=0.68008\bit ,
        \]
        所以 $\I{X;Z}=1-0.68008=0.31992\bit$。
  \item 结论：$0.31992<0.53101$，确实变少了。
        （顺便得到一个规律：只要原来的错误率小于 $\frac12$，再独立地错一次
        总比原来错得更多，所以级联一次比一次糟——第二次纯属浪费胶卷。）
\end{itemize}
\end{calc}

\begin{qa}{“信息处理只会减少信息”——会不会有例外？}
\textbf{会——只要 $X\to Y\to Z$ 不是马尔可夫链，这个不等式就可能不成立。}
用的时候两条前提缺一不可：（1）$X\to Y\to Z$ 必须真的是条链；
（2）比的是“$\I{X;Z}$ 与 $\I{X;Y}$”这种	extbf{同一方向}的量，
不是拿条件互信息去比。
最著名的反例就是前面链式法则里的\textbf{例 B（异或）}：$X,Y$ 独立，
$\I{X;Y}=0$；但加上 $Z=X\oplus Y$ 之后，$\I{X;Z|Y}=1$ bit。
也就是说
\[
\I{X;Y}=0\ \ <\ \ \I{X;Z|Y}=1 .
\]
人话：\textbf{“条件化”可以让互信息变大（$0\to1$），也可以让它变小}；
但\textbf{顺着马尔可夫链往前传}（$X\to Y\to Z$ 这种方向）只会变小。
这两个方向千万别混——这是本章最容易丢分的一句话。
\end{qa}

% ==================================================================
\begin{exam}{考试怎么考（第 3 章只考这四类题）}
\textbf{题型 1：给联合分布表，求 $\I{X;Y}$。}（最高频，几乎必有一道）

\emph{答题就写这 5 行，一行都别省：}
\begin{enumerate}[leftmargin=1.8em,itemsep=1pt,topsep=1pt]
  \item 抄表并算出边缘分布 $p(x),p(y)$（并把两者加起来各验一次 $=1$）；
  \item 写出 $\Hx{X}$、$\Hx{Y}$、$\Hx{X,Y}$ 的数值（或写成 $-\sum p\lb p$ 带数字）；
  \item 写出所用公式，如 $\I{X;Y}=\Hx{X}+\Hx{Y}-\Hx{X,Y}$；
  \item 代入数字，写出结果与单位（\textbf{bit}）；
  \item 校验：$0\le\I{X;Y}\le\min\{\Hx{X},\Hx{Y}\}$，若超出范围立刻回头查表。
\end{enumerate}

\textbf{题型 2：证明熵与互信息的关系式}（如 $\Hx{X,Y,Z}=\Hx{X}+\Hx{Y|X}+\Hx{Z|XY}$、
$\I{X;YZ}=\I{X;Y}+\I{X;Z|Y}$ 之类）。\emph{套路固定}：先把所有 $\I{\cdot}$ 用
$\Hx{a}-\Hx{a|b}$ 或 $\Hx{a}+\Hx{b}-\Hx{a,b}$ 展开成 $H$ 的加减，
\textbf{加一项减一项凑定义}，让中间项相互抵消，最后只剩非负量。
\textbf{关键：每一步都要注明依据}（“由链式法则”“由条件熵非负”），
阅卷按步给分。

\textbf{题型 3：判断独立 / 求极值 / 判断某不等式成立。}
独立 $\iff\I{X;Y}=0\iff p(x,y)=p(x)p(y)$ 处处成立；
$\I{X;Y}$ 的最大值 $\min\{\Hx{X},\Hx{Y}\}$ 在“一方完全决定另一方”时取到。

\textbf{题型 4：数据处理定理 $X\to Y\to Z\Rightarrow\I{X;Z}\le\I{X;Y}$。}
照抄前面数据处理那三道工序，第 2 步一定要\textbf{写出马尔可夫条件的表达式}
$p(z|xy)=p(z|y)$，再说明“故 $\I{X;Z|Y}=0$”。

\medskip
\noindent\textbf{易错点（每一条都有人栽过）}
\begin{itemize}[leftmargin=1.6em,itemsep=1pt,topsep=1pt]
  \item $\Hx{X|Y}\ne\Hx{Y|X}$（一般不等，例 1 里是 $0.875$ 对 $0.846$），
        但 $\I{X;Y}=\I{Y;X}$（严格相等）。\textbf{对称的只有互信息，不是条件熵。}
  \item 别把 $\I{X;Y}$ 和 $\Hx{X|Y}$ 搞混：前者是两圆的\textbf{重叠}，
        后者是左圆\textbf{露出来的月牙}。方向也相反：
        $\I{X;Y}$ 越大，$\Hx{X|Y}$ 越小（$\I{X;Y}=\Hx{X}-\Hx{X|Y}$）。
  \item $\I{x;y}$（小写，单个事件）可正可负；$\I{X;Y}$（大写，平均）
        \textbf{恒非负}。算出负数必是算错，别顺着往下写。
  \item 分子写联合概率、分母写两个边缘概率的\textbf{乘积}；
        别写成 $\lb\frac{p(x|y)}{p(y)}$ 这种把分母漏掉一项的形式。
  \item 权重用 $p(x,y)$；忘了权重（直接把四个事件互信息加起来）是最常见的错。
  \item 单位是 bit（底 $2$）；题目若用 $\ln$，单位是 nat，别混。
  \item 单看某一格 $\I{x;y}$ 可能是负数，这不矛盾（例 1 里就有 $-1$ 和 $-0.585$）；
        但平均值 $\I{X;Y}$ 一定 $\ge0$，两类量别混。
  \item 要证独立，必须 $p(x,y)=p(x)p(y)$ \textbf{每一格}都成立（或直接算出
        $\I{X;Y}=0$）；只在某一格上相等说明不了什么。
\end{itemize}
\end{exam}

% ==================================================================
\begin{qa}{你可能想问}
\textbf{问 1：互信息和熵到底什么关系？可以互相算出来吗？}

答：\textbf{不能互相算，它们是“两个圆”的关系}。
$\Hx{X}$、$\Hx{Y}$ 是两个圆的面积，$\I{X;Y}$ 是重叠面积。
知道两个圆和一个重叠，还差一个“并集”才能确定全部，所以单靠 $\Hx{X}$ 和
$\Hx{Y}$ 定不出 $\I{X;Y}$：前面手算过的天气表和“独立表”$(0.3,0.2,0.3,0.2)$ 的
边缘分布\textbf{完全一样}（都是 $0.5/0.5$ 和 $0.6/0.4$），
但一个是 $0.12451$ bit，一个是 $0$ bit。
记住三句：$\I{X;X}=\Hx{X}$（完全相关时重叠 $=$ 全圆）、
独立时 $\I{X;Y}=0$（不重叠）、其余介于两者之间。
真要靠 $H$ 算 $I$，必须\textbf{再多知道一个} $H(X,Y)$（或 $H(X|Y)$、$H(Y|X)$ 之一）。

\textbf{问 2：为什么互信息不会超过 $\Hx{X}$？会不会出现 $\I{X;Y}>\Hx{X}$？}

答：看第一条公式 $\I{X;Y}=\Hx{X}-\Hx{X|Y}$。$\Hx{X|Y}$ 是一堆 $-\lb p(x|y)$
的加权平均，每个 $p(x|y)\le1$，所以 $-\lb p(x|y)\ge0$，平均当然 $\ge0$。
\textbf{“剩下的不确定性”不可能是负的}——你不可能“知道得比全部还多”。
极端情形 $\Hx{X|Y}=0$（知道 $Y$ 就等于知道 $X$）时取等 $\I{X;Y}=\Hx{X}$，
如“下雨一定带伞、不下雨一定不带”。同理 $\I{X;Y}\le\Hx{Y}$，
所以上限是 $\min\{\Hx{X},\Hx{Y}\}$。

\textbf{问 3：为什么不直接用 $p(x,y)$ 的大小判断相关性，非要绕一个比值加 $\lb$？}

答：因为 $p(x,y)$ 的数值\textbf{同时}被“$X$ 有多常见”和“$Y$ 有多常见”压着，
横向比较会误判。看本讲两个例子：
例 1 里 $p(r,u)=0.4$（挺大一个数），对应的互信息只有 $0.415$ bit；
例 3 里 $p(x=6,y=12)=\frac1{36}=0.0278$（小得可怜），对应的互信息却是
$2.585$ bit。\textbf{要是看概率大小，就会得出“天气和带伞比骰子更相关”
这种荒唐结论。}
比值 $\dfrac{p(x,y)}{p(x)p(y)}$ 把“各自多常见”除掉了（分母就是“没关系时
本该有的概率”），于是所有关联都被折算到同一把尺子上：
比值 $=1$ 就是“毫无关联”这条零线，$>1$ 正相关，$<1$ 负相关。
取 $\lb$ 则是为了让“独立 $\Rightarrow$ 相加”：独立时比值恒为 $1$，
$\lb1=0$，信息量正好可以一节一节往上加（前面解释对数时讲过）。
一句话：\textbf{比值负责“除掉常见程度的干扰”，$\lb$ 负责“把倍数翻译成可加的 bit”。}
\end{qa}
